\chapter{Probabilistic Foundations and Limit Theorems}

\section{Probability and Moment Bounds}

\subsection{Foundations of Probability Theory}

To discuss limit theorems rigorously, we must first establish the measure-theoretic foundation of probability. This framework prevents the paradoxes that arise in naive probability and allows us to handle continuous and infinite-dimensional spaces seamlessly.

\begin{definition}[Probability Space]
A \textit{probability space} is a triple $(\Omega, \mathcal{F}, \Prob)$ consisting of:
\begin{enumerate}
\item \textit{Sample space ($\Omega$)}: a nonempty set of all possible outcomes of a random experiment. Each element $\omega$ represents one individual outcome.
\item \textit{$\sigma$-algebra of events ($\mathcal{F}$)}: a collection of subsets of $\Omega$, called events, that contains $\Omega$ and is closed under complements relative to $\Omega$ and countable unions. Explicitly, for every event $A$ and every sequence of events $A_n$, these conditions are
\[
\begin{gathered}
\Omega \in \mathcal{F}, \\
\Omega \setminus A \in \mathcal{F}, \\
\bigcup_{n=1}^{\infty} A_n \in \mathcal{F}.
\end{gathered}
\]
\item \textit{Probability measure ($\Prob$)}: a function assigning to each event a number between zero and one, with total probability one:
\[
\begin{gathered}
\Prob: \mathcal{F} \to [0,1], \\
\Prob(\Omega) = 1.
\end{gathered}
\]
It is countably additive: for every sequence of pairwise disjoint events $A_n$, the probability of their union equals the sum of their probabilities:
\[
\Prob\!\left(\bigcup_{n=1}^{\infty} A_n\right)
= \sum_{n=1}^{\infty} \Prob(A_n).
\]
\end{enumerate}
\end{definition}

\begin{remark}[Why a $\sigma$-algebra?]
We cannot simply assign probabilities to \textit{all} subsets of $\Omega$ (for example, if $\Omega = \R$, the Axiom of Choice implies the existence of non-measurable sets like the Vitali set). The $\sigma$-algebra $\mathcal{F}$ restricts the domain of $\Prob$ to a collection of ``well-behaved'' sets. Closure under countable unions is strictly necessary because probability theory frequently involves limits of events (e.g., ``the event happens eventually''), and we must guarantee that the limit of a sequence of valid events is still a valid event to which we can assign a probability.
\end{remark}

\begin{definition}[Random Variable and Expectation]
A \textit{random variable} $S$ is a measurable function $S: (\Omega, \mathcal{F}) \to (\R, \mathcal{B}(\R))$, where $\mathcal{B}(\R)$ is the Borel $\sigma$-algebra on $\R$. Measurability means that for any Borel set $B$, the preimage $S^{-1}(B) = \{\omega \in \Omega : S(\omega) \in B\}$ is an element of $\mathcal{F}$.
For the purposes of this chapter, we restrict attention to integrable
random variables, meaning that
\[
\E[|S|]=\int_\Omega |S(\omega)|\,\dd\Prob(\omega)<\infty.
\]
For such a random variable, the \textit{expectation} (or expected value)
is the Lebesgue integral
\begin{equation}
\E[S] = \int_\Omega S(\omega) \dd \Prob(\omega)
\end{equation}
and is finite. More generally, extended expectations can be defined from
the positive and negative parts when they are not both infinite.
\end{definition}

\begin{remark}[Why Measurability?]
We require $S$ to be measurable so that we can calculate probabilities of events defined by $S$. For instance, to compute $\Prob(S \leq x)$, we need the set $\{\omega : S(\omega) \leq x\}$ to be in $\mathcal{F}$ so that the measure $\Prob$ can be applied to it. Without measurability, the expectation integral would be undefined.
\end{remark}

\input{figures/conditional-events.tex}

\subsection{Markov and Chebyshev Inequalities}

These inequalities provide deterministic bounds on the tails of probability distributions using only the moments (expectation and variance) of the random variable.

\begin{proposition}[Markov's Inequality]
Let $S$ be a nonnegative random variable ($S \geq 0$ almost surely) and let $a > 0$. Then:
\begin{equation}
\Prob(S \geq a) \leq \frac{\E[S]}{a} % <--- QUI USA \Prob
\end{equation}

\end{proposition}

\begin{proof}
We introduce the \textit{indicator function} $\mathbf{1}_{S \geq a}(\omega)$, which equals $1$ if $S(\omega) \geq a$ and $0$ otherwise.
Observe the pointwise inequality for all $\omega \in \Omega$:
\begin{equation}
S(\omega) \geq a \cdot \mathbf{1}_{S \geq a}(\omega)
\end{equation}

This holds because if $S(\omega) \geq a$, the right side is $a$, and $S(\omega) \geq a$; if $S(\omega) < a$, the right side is $0$, and $S(\omega) \geq 0$ by assumption.
Since the Lebesgue integral preserves inequalities for nonnegative functions, we integrate both sides with respect to $\Prob$ over $\Omega$:
\begin{equation}
\int_\Omega S(\omega) \dd \Prob(\omega) \geq \int_\Omega a \cdot \mathbf{1}_{S \geq a}(\omega) \dd \Prob(\omega)
\end{equation}

The left side is exactly $\E[S]$. On the right side, $a$ is a constant and can be pulled out. The integral of an indicator function over $\Omega$ is precisely the probability of the set it indicates:
\begin{equation}
\begin{aligned}
\E[S] &\geq a \int_\Omega \mathbf{1}_{S \geq a}(\omega) \dd \Prob(\omega) \\
&= a \cdot \Prob(S \geq a)
\end{aligned}
\end{equation}
Dividing both sides by $a > 0$ yields the result.
\end{proof}

\begin{definition}[Variance]
The \textit{variance} of a random variable $S$ with finite expectation $\mu = \E[S]$ is defined as the expected squared deviation from the mean:
\begin{equation}
\var(S) = \E[(S - \E[S])^2]
\end{equation}

\end{definition}

\begin{theorem}[Chebyshev's Inequality]
If $S$ has finite expectation $\mu$ and finite variance $\sigma^2 = \var(S) < \infty$, then for any $a > 0$:
\begin{equation}
\Prob(|S - \mu| \geq a) \leq \frac{\sigma^2}{a^2}
\end{equation}

\end{theorem}

\begin{proof}
Define a new random variable $X = (S - \mu)^2$. Since it is a square, $X \geq 0$ almost surely.
Notice that the event $\{|S - \mu| \geq a\}$ is mathematically equivalent to the event $\{(S - \mu)^2 \geq a^2\}$, which is exactly $\{X \geq a^2\}$.
We apply Markov's Inequality to the nonnegative random variable $X$ with the threshold $a^2 > 0$:
\begin{equation}
\Prob(X \geq a^2) \leq \frac{\E[X]}{a^2}
\end{equation}
Substituting back $X = (S - \mu)^2$ and recognizing that
\[
\begin{aligned}
\E[X] &= \E[(S - \mu)^2] \\
&= \var(S) \\
&= \sigma^2,
\end{aligned}
\]
we obtain:
\begin{equation}
\Prob(|S - \mu| \geq a) \leq \frac{\sigma^2}{a^2}
\end{equation}

\end{proof}

\section{Limit Theorems for Sums of Random Variables}

\subsection{Weak Law of Large Numbers}

The Weak Law of Large Numbers (WLLN) formalizes the intuitive idea that the average of a large number of independent observations converges to the expected value.

\begin{definition}[Independent and Identically Distributed (i.i.d.)]
A sequence of random variables $S_1, S_2, \dots$ is:
\begin{itemize}
\item \textit{Identically distributed}: if they all share the same cumulative distribution function, and consequently, the same expectation $\mu$ and variance $\sigma^2$.
\item \textit{Independent}: if the joint distribution of any finite subset is the product of their marginal distributions. Crucially for expectations, this implies that for $i \neq j$, the expectation of the product factors:

\[
\E[S_i S_j] = \E[S_i]\E[S_j].
\]
Consequently, their covariance is zero:
\[
\E[(S_i - \mu)(S_j - \mu)] = 0.
\]

\end{itemize}
\end{definition}

\begin{theorem}[Weak Law of Large Numbers]
Let $S_1, S_2, \dots$ be a sequence of i.i.d. random variables with mean $\mu$ and finite variance $\sigma^2$. Let the sample mean be
\[
z_m = (1/m) \sum_{i=1}^m S_i.
\]
Then for any $\varepsilon > 0$:
\begin{equation}
\lim_{m \to \infty} \Prob(|z_m - \mu| \geq \varepsilon) = 0
\end{equation}

\end{theorem}

\begin{proof}
We first explicitly calculate the expectation and variance of $z_m$.
By the linearity of the expectation operator:
\begin{equation}
\begin{aligned}
\E[z_m] &= \E\left[ \frac{1}{m} \sum_{i=1}^m S_i \right] \\
&= \frac{1}{m} \sum_{i=1}^m \E[S_i] \\
&= \frac{1}{m} (m \mu) \\
&= \mu
\end{aligned}
\end{equation}
Next, we calculate the variance of $z_m$. By definition:
\begin{equation}
\begin{aligned}
\var(z_m) &= \E[(z_m - \E[z_m])^2] \\
&= \E\left[ \left( \frac{1}{m} \sum_{i=1}^m (S_i - \mu) \right)^2 \right]
\end{aligned}
\end{equation}
Expanding the square of the sum yields a double sum:
\begin{equation}
\var(z_m) = \frac{1}{m^2} \E\left[ \sum_{i=1}^m \sum_{j=1}^m (S_i - \mu)(S_j - \mu) \right]
\end{equation}
Using the linearity of expectation again, we bring the expectation inside the sums:
\begin{equation}
\var(z_m) = \frac{1}{m^2} \sum_{i=1}^m \sum_{j=1}^m \E[(S_i - \mu)(S_j - \mu)]
\end{equation}
We evaluate the terms in the double sum based on the indices:
\begin{itemize}
\item \textit{Diagonal terms}: if $i = j$, the term is

\[
\begin{aligned}
\E[(S_i - \mu)^2] &= \var(S_i) \\
&= \sigma^2.
\end{aligned}
\]
There are exactly $m$ such terms on the diagonal of the double sum.
\item \textit{Off-diagonal terms}: if $i \neq j$, the term is $\E[(S_i - \mu)(S_j - \mu)]$. Because $S_i$ and $S_j$ are independent, this factors into

\[
\begin{aligned}
\E[S_i - \mu]\E[S_j - \mu] &= 0 \cdot 0 \\
&= 0.
\end{aligned}
\]
All $m^2 - m$ off-diagonal terms vanish.
\end{itemize}
Summing these up, we get:
\begin{equation}
\begin{aligned}
\var(z_m) &= \frac{1}{m^2} (m \sigma^2) \\
&= \frac{\sigma^2}{m}
\end{aligned}
\end{equation}
Now, we apply Chebyshev's Inequality to the random variable $z_m$ with threshold $\varepsilon > 0$:
\begin{equation}
\begin{aligned}
\Prob(|z_m - \mu| \geq \varepsilon) &\leq \frac{\var(z_m)}{\varepsilon^2} \\
&= \frac{\sigma^2}{m \varepsilon^2}
\end{aligned}
\end{equation}
Since $\sigma^2$ and $\varepsilon^2$ are fixed finite constants, taking the limit as $m \to \infty$ gives:
\begin{equation}
\lim_{m \to \infty} \frac{\sigma^2}{m \varepsilon^2} = 0
\end{equation}
Because probabilities are nonnegative, we conclude that
\[
\lim_{m \to \infty} \Prob(|z_m - \mu| \geq \varepsilon) = 0.
\]

\end{proof}

\subsection{Characteristic Functions and the Central Limit Theorem}


While the WLLN tells us the sample mean converges to a constant, the Central Limit Theorem (CLT) describes the exact distribution of the fluctuations around that constant.

\begin{definition}[Convergence in Distribution]
A sequence of random variables $Y_m$ \textit{converges in distribution} to a random variable $Y$, denoted $Y_m \xrightarrow{d} Y$, if their cumulative distribution functions (CDFs) $F_{Y_m}(y)$ converge to $F_Y(y)$ at every point $y$ where $F_Y(y)$ is continuous.
\end{definition}

\begin{remark}[Why use Characteristic Functions?]
Proving convergence of CDFs directly is analytically intractable. Instead, we use the \textit{Characteristic Function} (CF), which is the Fourier transform of the probability measure. The CF always exists, uniquely determines the probability distribution, and converts the convolution of independent random variables into simple algebraic multiplication, making limit calculations straightforward.
\end{remark}

\begin{definition}[Characteristic Function]
The characteristic function of a random variable $S$ is defined as:
\begin{equation}
\varphi_S(t) = \E[e^{itS}] = \int_\Omega e^{itS(\omega)} \dd \Prob(\omega), \quad t \in \R
\end{equation}

\end{definition}

\begin{theorem}[Lindeberg--Lévy Central Limit Theorem]
Let $S_1, S_2, \dots$ be a sequence of i.i.d. random variables with mean $\mu$ and finite variance $\sigma^2 > 0$. Let
\[
z_m = (1/m) \sum_{i=1}^m S_i.
\]

Then the normalized sample mean converges in distribution to a standard normal random variable:
\begin{equation}
Y_m = \frac{\sqrt{m}(z_m - \mu)}{\sigma} \xrightarrow{d} \mathcal{N}(0, 1)
\end{equation}

\end{theorem}

\subsection{Proof of the Central Limit Theorem}

\begin{proof}
We will prove this by showing that the characteristic function of $Y_m$ converges pointwise to the characteristic function of the standard normal distribution $\mathcal{N}(0,1)$.

\textit{Step 1: Express $Y_m$ as a sum of centered variables.}
Let $X_i = S_i - \mu$. The variables $X_i$ are i.i.d. with $\E[X_i] = 0$ and $\var(X_i) = \sigma^2$. We can rewrite $Y_m$ as:
\begin{equation}
\begin{aligned}
Y_m &= \frac{\sqrt{m}}{\sigma} \left( \frac{1}{m} \sum_{i=1}^m X_i \right) \\
&= \frac{1}{\sigma\sqrt{m}} \sum_{i=1}^m X_i
\end{aligned}
\end{equation}

\textit{Step 2: Relate the CF of $Y_m$ to the CF of $X_1$.}
Using the properties of characteristic functions:
\begin{itemize}
\item \textit{Scaling}: the following identity holds:

\[
\begin{aligned}
\varphi_{cX}(t) &= \E[e^{it(cX)}] \\
&= \varphi_X(ct).
\end{aligned}
\]

\item \textit{Independence}: if $X$ and $Y$ are independent,

\[
\begin{aligned}
\varphi_{X+Y}(t) &= \E[e^{itX}e^{itY}] \\
&= \E[e^{itX}]\E[e^{itY}] \\
&= \varphi_X(t)\varphi_Y(t).
\end{aligned}
\]

\end{itemize}
Applying these properties to $Y_m$:
\begin{equation}
\begin{aligned}
\varphi_{Y_m}(t) &= \varphi_{\sum X_i / (\sigma\sqrt{m})}(t) \\
&= \left[ \varphi_{X_1}\left( \frac{t}{\sigma\sqrt{m}} \right) \right]^m
\end{aligned}
\end{equation}

\textit{Step 3: Taylor expand the characteristic function of $X_1$.}
We expand $\varphi_{X_1}(u)$ around $u=0$. To do this rigorously, we differentiate under the integral sign. This is justified by the Dominated Convergence Theorem because the derivatives of the integrand are bounded by $|X_1|^k$, and we assume $\E[|X_1|^2] < \infty$.
\begin{align*}
\varphi_{X_1}(u) &= \E[e^{iuX_1}] \\
    \varphi_{X_1}'(u) &= \E[iX_1 e^{iuX_1}] \implies \varphi_{X_1}'(0) = i\E[X_1] = 0 \\
    \varphi_{X_1}''(u) &= \E[(iX_1)^2 e^{iuX_1}] = -\E[X_1^2 e^{iuX_1}] \implies \varphi_{X_1}''(0) = -\E[X_1^2] = -\var(X_1) = -\sigma^2
\end{align*}
By Taylor's theorem with the Peano form of the remainder, we have:
\begin{equation}
\begin{aligned}
\varphi_{X_1}(u) &= \varphi_{X_1}(0) + u\varphi_{X_1}'(0) + \frac{u^2}{2}\varphi_{X_1}''(0) + o(u^2) \\
&= 1 - \frac{\sigma^2 u^2}{2} + o(u^2).
\end{aligned}
\end{equation}
Here
\[
o(u^2)
\]
 denotes a function such that
\[
\lim_{u \to 0} o(u^2)/u^2 = 0.
\]

\textit{Step 4: Substitute and take the limit.}
We substitute $u = (t/(\sigma\sqrt{m}))$ into the Taylor expansion. Note that as $m \to \infty$, $u \to 0$:
\begin{equation}
\begin{aligned}
\varphi_{X_1}\left( \frac{t}{\sigma\sqrt{m}} \right) &= 1 - \frac{\sigma^2}{2} \left( \frac{t}{\sigma\sqrt{m}} \right)^2 + o\left( \frac{t^2}{\sigma^2 m} \right) \\
&= 1 - \frac{t^2}{2m} + o\left( \frac{1}{m} \right)
\end{aligned}
\end{equation}
Now, we substitute this back into the expression for $\varphi_{Y_m}(t)$:
\begin{equation}
\varphi_{Y_m}(t) = \left[ 1 - \frac{t^2}{2m} + o\left( \frac{1}{m} \right) \right]^m
\end{equation}

To evaluate the limit as $m \to \infty$, we rewrite the expression using the exponential and natural logarithm:
\begin{equation}
\varphi_{Y_m}(t) = \exp\left( m \ln\left( 1 - \frac{t^2}{2m} + o\left( \frac{1}{m} \right) \right) \right)
\end{equation}
Using the Taylor expansion
\[
\ln(1 + x) = x + o(x)
\]
 for small $x$, where
\[
x = -(t^2/(2m)) + o( (1/m) ):
\]

\begin{equation}
\begin{aligned}
m \ln\left( 1 - \frac{t^2}{2m} + o\left( \frac{1}{m} \right) \right) &= m \left( -\frac{t^2}{2m} + o\left( \frac{1}{m} \right) \right) \\
&= -\frac{t^2}{2} + m \cdot o\left( \frac{1}{m} \right)
\end{aligned}
\end{equation}
The remainder term satisfies
\[
\lim_{m \to \infty} m \cdot o(1/m) = 0,
\]
Therefore, the exponent converges to $-t^2/2$. Therefore:
\begin{equation}
\lim_{m \to \infty} \varphi_{Y_m}(t) = e^{-t^2/2}
\end{equation}

\textit{Step 5: Apply Lévy's Continuity Theorem.}
\textit{Lévy's Continuity Theorem} states that if a sequence of characteristic functions $\varphi_{Y_m}(t)$ converges pointwise to a limit function $\varphi(t)$ that is continuous at $t=0$, then $Y_m$ converges in distribution to a random variable $Y$ whose characteristic function is $\varphi(t)$.
The limit function $e^{-t^2/2}$ is continuous everywhere. Furthermore, it is exactly the characteristic function of the standard normal distribution $\mathcal{N}(0,1)$.
Thus, we conclude that $Y_m \xrightarrow{d} \mathcal{N}(0, 1)$. \end{proof}

\subsection{Normal Approximation and Interpretation}

Figure~\ref{fig:ch16-clt} illustrates the approach to a normal distribution described by the Central Limit Theorem:
\[
z_m \xrightarrow{d} N(0,1).
\]

\begin{figure}[!htbp]
\centering
\begin{tikzpicture}
\begin{axis}[
  width=13cm,height=8cm, axis lines=left,
  xlabel={$z=((\sqrt{m}(\bar X_m-\mu))/\sigma)$}, ylabel={density},
  xmin=-4,xmax=4,ymin=0,ymax=0.62,
  legend pos=north east, legend style={font=\scriptsize, draw=black!30}
  ]
  \addplot[black!45,thick,const plot,domain=-4:4,samples=9]
    coordinates {(-3.5,0.05)(-2.5,0.09)(-1.5,0.22)(-0.5,0.4)(0.5,0.4)(1.5,0.22)(2.5,0.09)(3.5,0.05)};
  \addlegendentry{histogram of $z_m$, moderate $m$}
  \addplot[figred,very thick,domain=-4:4,samples=150] {1/sqrt(2*pi)*exp(-x^2/2)};
  \addlegendentry{$N(0,1)$ limit}
\end{axis}
\end{tikzpicture}
\caption[Central Limit Theorem]{Illustrative histogram of a standardized sample mean for a moderate sample size, compared with the standard normal limiting density (red). The horizontal coordinate is $z=\sqrt{m}(\bar X_m-\mu)/\sigma$.}
\label{fig:ch16-clt}
\end{figure}
